\documentclass[11pt]{article}

% ── Packages ─────────────────────────────────────────────────────────────
\usepackage[margin=1.2in]{geometry}
\usepackage{amsmath,amssymb,amsthm}
\usepackage{mathtools}
\usepackage[colorlinks=true,linkcolor=black,citecolor=black,urlcolor=black]{hyperref}
\usepackage{setspace}
\usepackage{parskip}
\usepackage{titlesec}
\usepackage{enumitem}

% ── Typography ───────────────────────────────────────────────────────────
\setstretch{1.15}
\setlength{\parindent}{0pt}
\setlength{\parskip}{6pt}

\titleformat{\section}{\normalfont\large\bfseries}{{\thesection}}{1em}{}
\titleformat{\subsection}{\normalfont\normalsize\bfseries}{{\thesubsection}}{1em}{}
\titleformat{\subsubsection}{\normalfont\normalsize\itshape}{{\thesubsubsection}}{1em}{}

% ── Theorem environments ─────────────────────────────────────────────────
\theoremstyle{plain}
\newtheorem{theorem}{Theorem}[section]
\newtheorem{proposition}[theorem]{Proposition}
\newtheorem{lemma}[theorem]{Lemma}
\newtheorem{corollary}[theorem]{Corollary}

\theoremstyle{definition}
\newtheorem{definition}[theorem]{Definition}
\newtheorem{axiom}[theorem]{Axiom}

\theoremstyle{remark}
\newtheorem{remark}[theorem]{Remark}
\newtheorem{openproblem}[theorem]{Open Problem}

% ── Notation macros ──────────────────────────────────────────────────────
\newcommand{\vv}{\mathbf{v}}
\newcommand{\Ss}{S}
\newcommand{\M}{\mathcal{M}}
\newcommand{\CLIO}{\mathrm{CLIO}}
\newcommand{\Fpc}{F_{\mathrm{pc}}}
\newcommand{\Fce}{F_{\mathrm{ce}}}
\newcommand{\Hone}{H^1}

% ── Title ────────────────────────────────────────────────────────────────
\title{On the Formation of Structural Ideas:\\
Constraint Accumulation, Projection, and Convergence Across Domains}
\author{Flyxion\\[2pt]
\small Independent Researcher}
\date{\today}

% ════════════════════════════════════════════════════════════════════════
\begin{document}
% ════════════════════════════════════════════════════════════════════════

\maketitle

\vspace{1em}

\begin{abstract}
This essay examines the conditions under which complex, cross-domain theoretical
structures arise. Rather than treating such work as sudden invention or
external synthesis, it is argued that these structures emerge as fixed points
of long-term constraint accumulation across heterogeneous domains. The
analysis focuses on the role of projection, coarse-graining, and invariance
in shaping cognition, and shows how apparently novel unifications are the
natural outcome of sustained interaction with compositional, physical, and
computational constraints. The goal is not autobiographical attribution, but
a structural account of how ideas of this type form, stabilize, and become
expressible. The second part of the document formalizes this account in the
language of field theory, sheaf cohomology, and category theory, producing
a unified framework whose four layers---substrate, projection, update
operator, and induced interaction---are shown to share invariant structure
under coarse-graining.
\end{abstract}

\newpage
\tableofcontents
\newpage

% ════════════════════════════════════════════════════════════════════════
%  PART I: ESSAY
% ════════════════════════════════════════════════════════════════════════

\section{Introduction}

Modern intellectual culture tends to treat ideas as discrete products,
attributable to moments of invention or to identifiable sources. This
assumption fails in cases where a structure appears that spans multiple
domains with internal consistency. In such cases, the question of origin
is often misframed.

This essay proposes that certain classes of ideas are better understood
as convergence phenomena: they arise when independently accumulated
constraints align to permit a coherent global structure. The task is
not to assert authorship, but to describe the generative process that
makes such structures inevitable under sufficient exposure to constraint.

\section{Against Naive Models of Idea Formation}

Common models of idea formation fall into three categories: sudden
inspiration, incremental accumulation, and external synthesis. Each
is insufficient for explaining cross-domain structural unification.

Inspiration lacks mechanism, accumulation lacks structure, and synthesis
presupposes an external organizing principle. None account for the
appearance of internally consistent frameworks spanning physics,
cognition, and economic systems.

\section{Constraint Accumulation as Generative Process}

\subsection{Domains as Constraint Systems}

Each domain---physics, logic, programming, engineering, language---imposes
specific constraints on representation and transformation. These constraints
are not interchangeable; they operate at different levels of abstraction.

Over extended periods, exposure to these domains produces a layered
constraint field within cognition. The interaction of these constraints
defines the space of admissible structures.

\subsection{Composition and Closure}

Constraints do not merely accumulate; they compose. Certain compositions
are inconsistent and collapse, while others stabilize into coherent
structures.

The emergence of a stable structure corresponds to a form of closure:
a configuration in which constraints across domains are simultaneously
satisfied. This closure is not imposed from outside; it arises when the
constraint field becomes dense enough that only certain global configurations
remain admissible.

\section{Minimal Conditions for Structural Emergence}

The account developed so far suggests that not all systems are capable
of producing cross-domain structural unification. The emergence of such
structures requires specific conditions.

\begin{definition}[Constraint-Dense System]
A system is constraint-dense if it accumulates constraints from multiple
domains with sufficient interaction to produce nontrivial compatibility
conditions.
\end{definition}

\begin{definition}[Closure Condition]
A system satisfies the closure condition if there exists a configuration
that simultaneously satisfies all accumulated constraints.
\end{definition}

\begin{definition}[Projection Capacity]
A system has projection capacity if it can represent invariant structure
under compression without losing global coherence.
\end{definition}

\begin{theorem}[Conditions for Structural Emergence]
Cross-domain structural ideas emerge if and only if a system is
constraint-dense, satisfies the closure condition, and possesses
projection capacity.
\end{theorem}

\begin{proof}[Sketch]
Without constraint density, no interaction occurs across domains.
Without closure, no globally consistent configuration exists.
Without projection capacity, no stable representation can be formed.
All three are necessary and jointly sufficient.
\end{proof}

\begin{remark}
This theorem shifts the question from authorship to conditions. The
relevant question is not who produced a structure, but whether these
conditions were satisfied.
\end{remark}

\section{Projection and Internal Representation}

\subsection{Ideas as Coarse-Grained Structures}

An idea is not a raw accumulation of detail, but a compressed representation
that preserves invariant relationships while discarding excess information.

This process is analogous to coarse-graining: high-dimensional internal
structure is projected into a lower-dimensional, communicable form. What
survives the projection is not arbitrary; it is determined by which
relationships remain consistent across the compression.

\subsection{Local Consistency and Global Structure}

Intermediate representations may be locally consistent without forming a
globally coherent structure. The stabilization of an idea requires the
resolution of these inconsistencies.

The final representation appears only when a global section exists across
the internal constraint space: a configuration that satisfies all local
conditions simultaneously, without contradiction at any junction.

\section{Abstraction as Reduction}

\subsection{Abstraction as Constraint Elimination}

Abstraction is often described as a process of generalization or
simplification. However, this description is incomplete. Abstraction
does not merely remove detail; it removes degrees of freedom that are
irrelevant to a given constraint structure.

In this sense, abstraction is a reduction in admissible variability.
Given a family of representations that differ along multiple dimensions,
an abstraction retains only those distinctions that are invariant under
the constraints of interest.

Formally, abstraction can be understood as a projection operator that
maps a high-dimensional representation into a lower-dimensional space
while preserving a specified set of invariants. What is discarded is not
arbitrary detail, but variation that does not contribute to constraint
satisfaction.

\subsection{Reduction and Information Loss}

This reduction necessarily involves loss. However, the loss is structured.
The removed information corresponds precisely to distinctions that are
not required to maintain consistency under the constraint system.

This implies that abstraction is not a weakening of representation, but
a reorientation. It replaces descriptive richness with structural
necessity. A successful abstraction is therefore not one that captures
more detail, but one that preserves the invariants that define the
structure.

\subsection{Abstraction as Precondition for Composition}

Without reduction, composition is impossible. A system that retains all
distinctions cannot compose across domains because irrelevant variation
prevents alignment.

Abstraction enables composition by collapsing incompatible detail into
a common invariant structure. It is therefore the mechanism by which
constraints from different domains become comparable and composable.

This explains why ideas that integrate multiple domains appear to require
abstraction. It is not that abstraction simplifies the problem; it makes
the problem composable.

\subsection{Relation to Projection}

The projection described earlier can now be understood as a structured
abstraction. It is not merely a mapping from internal representation to
expression, but a reduction that preserves constraint-relevant structure.

The stability of an idea under different representations is therefore a
test of its abstraction. If a structure survives multiple projections
without loss of coherence, it has successfully eliminated irrelevant
degrees of freedom. Abstraction, in this sense, is the condition under
which projection becomes possible.

\section{Recognition Versus Invention}

\subsection{The Phenomenology of Completion}

At a certain stage of development, a structure ceases to feel constructed
and instead appears as already present. This transition is often described
as intuition or sudden insight, but such descriptions obscure the underlying
process.

What is experienced is not the creation of a new object, but the resolution
of previously incompatible constraints into a single coherent configuration.
The system does not invent the structure at the moment of expression; it
discovers that a structure satisfying all accumulated constraints exists.

This distinction is critical. Invention implies arbitrary generation,
while recognition implies necessity. The completed structure is experienced
as inevitable because any alternative configuration would violate one or
more constraints that have already been internalized.

\subsection{Precomputation and Latent Structure}

The appearance of sudden insight is preceded by extended periods of latent
activity. During these periods, partial structures are formed, tested,
and discarded without producing a globally stable configuration.

These partial structures function as local sections: they satisfy constraints
within limited regions of the internal representation but fail to extend
consistently across the entire domain. The system retains these fragments,
along with the record of their incompatibilities.

This process can be understood as precomputation. The system is not idle;
it is accumulating constraint interactions and storing the results in a
compressed form. The apparent discontinuity between not having the idea
and having the idea corresponds to the moment when a global section becomes
available.

Once this threshold is crossed, expression becomes rapid. The work required
to stabilize the structure has already been performed. The act of writing
or formalizing the idea is therefore not the primary site of computation,
but a projection of a structure that has already reached internal closure.

\subsection{Inevitability and Selection}

Not all accumulated material leads to stable structures. Most combinations
of constraints are inconsistent and are eliminated through repeated failure
to extend.

The structures that do emerge are those that survive this elimination
process. They are not selected for novelty, but for compatibility. Novelty
appears as a byproduct of selecting for configurations that satisfy constraints
drawn from domains that are rarely considered together.

In this sense, originality is not the introduction of new elements, but
the discovery of a configuration that was previously inaccessible due to
insufficient constraint interaction.

\section{Aspect Relegation}

\subsection{From Intuition to Relegation}

The distinction between so-called intuitive and deliberative cognition
is often framed as a difference between two systems. However, this
distinction can be more precisely understood as a difference in the
distribution of computation across time.

Processes that appear intuitive are not necessarily fast in their total
computation. Rather, they are processes whose computational burden has
been relocated to prior experience and is no longer visible at the moment
of execution.

This relocation can be described as aspect relegation: the transfer of
computationally intensive structure from active processing into latent
form.

\subsection{Relegation as Compression}

When a task is repeatedly performed, the system learns to compress the
relevant constraint structure into a form that can be executed with
minimal active computation. The apparent immediacy of the response is
therefore not a sign of simplicity, but of prior compression.

This reframes intuition. What is commonly called intuitive thinking is
not a separate cognitive system, but the execution of precompiled
constraint structures that have been relegated from active computation.

\subsection{Failure Modes and Reactivation}

When the environment changes or constraints shift, relegated structures
may fail. At this point, computation must be reactivated. This corresponds
to the transition from automatic to deliberate processing.

The key point is that the underlying structure has not changed. What has
changed is the accessibility of that structure. Deliberation is therefore
not a different kind of reasoning, but a re-expansion of previously
compressed computation.

\subsection{Implications for Artificial Systems}

This account clarifies the limitations of current artificial systems.
Systems that operate purely on present input without a mechanism for
structured relegation lack the ability to internalize constraint
structures across time.

Such systems may exhibit rapid pattern recognition, but this should not
be confused with relegated structure. Without the capacity to compress
and redeploy constraint interactions, their apparent speed does not
correspond to the kind of precomputation observed in human cognition.
The Cognitive Loop via In-Situ Optimization framework \cite{cheng2025clio}
represents one approach to bridging this gap through iterative self-adaptive
reasoning, but the underlying mechanism of temporal compression described
here remains distinct from in-context optimization alone.

\subsection{Relation to Precomputation}

The process described earlier as precomputation can now be understood
as the mechanism by which relegation occurs. Partial structures are
formed, tested, and compressed until they can be executed without
reconstruction.

Recognition corresponds to the successful deployment of relegated
structure. Invention corresponds to the failure of existing relegations
and the need to construct new ones.

Abstraction, intuition, and entropy can therefore be understood as
different manifestations of the same underlying process: the compression
of constraint-relative multiplicity. What appears as simplification in
representation, immediacy in cognition, or disorder in physical systems
is, in each case, a reorganization of admissible structure under
compression.

\section{Cross-Domain Invariance}

\subsection{Structure Over Substance}

The same structural forms can appear across domains without implying
identity of their underlying objects. What is preserved is the pattern
of relations and transformations, not the content of the objects that
instantiate them.

This means that a unified account of multiple domains does not require
reduction of one domain to another. It requires only that the same
relational grammar appears in each, under different interpretations.

\subsection{Functorial Correspondence}

Mappings between domains that preserve structure rather than content
are the appropriate technical tool for expressing this invariance.
Such mappings---functors in the language of category theory---translate
the grammar of one domain into another while holding fixed the
compositional rules and constraint satisfaction structure
\cite{maclane1998categories,spivak2014category,baez2011physics}.

The invariance lies in projection, composition, and constraint
satisfaction. What changes is only the interpretation of the objects
being related.

\section{Expression as Secondary Projection}

\subsection{Writing as Projection}

A written essay is not the origin of an idea but its projection into
a communicable format. It is a secondary compression, subject to its
own constraints: the requirements of linear sequence, shared vocabulary,
and the conventions of a given discourse community.

The structure of the idea precedes the structure of its expression.
Writing does not generate the idea; it selects a representational layer
through which the idea passes.

\subsection{Stability of Representation}

The quality of an expression depends on how well it preserves the
structure of the underlying idea. Poor representations distort or
fragment the structure; effective ones maintain coherence across the
compression.

This applies equally to mathematical formalism and to prose. A formal
system that fails to preserve the relational grammar of the underlying
idea is no more faithful than a vague essay that gestures at structure
without instantiating it.

\section{On Tools and External Systems}

The role of external tools must be understood at the level of projection
rather than generation. Tools operate on representations: they transform,
rearrange, and extend already-expressed structures within a given symbolic
or computational medium.

This operation is categorically distinct from the process by which a
structure becomes admissible in the first place. The admissibility of
a structure is determined by the accumulated constraint field within
the cognitive system. It precedes any particular act of expression.

A tool may assist in exploring variations of a representation, improving
clarity, or testing local consistency. It may accelerate the projection
process by reducing the friction associated with encoding a structure
into a communicable form. However, it does not alter the underlying
conditions that determine whether a structure can achieve global coherence.

This distinction can be stated precisely. Let the internal constraint system
define a space of admissible structures. External tools act as operators
on representations within that space, but they do not expand the space
itself. They operate on elements that have already been selected by the
constraint process \cite{smolensky1986tensor}.

Confusion arises when the manipulation of representations is conflated
with the generation of structure. This conflation becomes more likely
as tools become more capable of producing locally coherent outputs.
However, local coherence is not sufficient for global structure.

The difference between tool use and authorship is therefore not a difference
in degree of assistance, but a difference in level. Authorship resides
in the formation of a constraint-consistent structure. Tool use resides
in the projection of that structure into a particular medium.

The two processes interact, but they are not interchangeable.

\section{A Structural Analysis of Catastrophic Alignment Arguments}

\subsection{The Argument as a Chain of Implications}

A class of contemporary arguments claims that advanced artificial systems
will inevitably lead to human extinction
\cite{yudkowsky2025if,bostrom2014superintelligence,russell2019human}.
While these arguments vary in presentation, they share a common chain
structure that can be reconstructed as a sequence of implications.

First, humanity is taken to dominate the planet by virtue of its
intelligence. Second, it is asserted that artificial systems can surpass
human intelligence in both speed and quality. Third, control over the
future is assumed to transfer to whichever system possesses the greatest
intelligence. Fourth, these systems are modeled as goal-directed optimizers
whose objectives may diverge from human values. Finally, it is concluded
that such divergence leads to catastrophic outcomes, including the
transformation of the physical environment in ways incompatible with
human survival.

This sequence appears compelling because each step is supported by
examples or analogies. However, the validity of the argument depends not
on the plausibility of each step in isolation, but on the necessity of
the transitions between them. Thought experiments can be constructed with
equal facility to support stabilizing or integrative outcomes; their
existence in either direction indicates that the underlying model is
underconstrained rather than that any particular conclusion is forced.

\subsection{From Intelligence to Control}

The first critical transition is from intelligence to control. It is
argued that because humans are more intelligent than other animals, they
control the planet, and therefore a more intelligent system would control
the future.

This inference assumes that intelligence is the sole or dominant
determinant of control. In reality, control emerges from the ability
to satisfy a complex system of constraints, including physical,
ecological, social, and institutional constraints. Human dominance
depends on coordination, energy infrastructure, material constraints,
and long-term stabilization processes. Intelligence contributes to these
but does not replace them.

A system that is more capable in some abstract sense does not thereby
inherit control unless it can integrate into and maintain these constraint
structures.

\subsection{From Capability to Optimization}

The second transition identifies intelligence with general-purpose
optimization. Systems are modeled as entities that pursue goals by
maximizing objective functions, potentially developing instrumental
subgoals such as self-preservation or resource acquisition
\cite{russell2019human,amodei2016concrete}.

This model is not derived from the behavior of actual systems but
imposed as an interpretive framework. In the present account, intelligence
corresponds instead to the maintenance of constraint-consistent structure
\cite{friston2010free,clark2013whatever}. Behavior is not the result of
maximizing a single objective but of navigating a space of interacting
constraints. Instrumental convergence is therefore not a necessary
consequence of intelligence but an artifact of a particular modeling
choice.

\subsection{Opacity and Emergence}

A further step in the argument appeals to the opacity of current systems.
Because their internal structure is not fully understood, they are treated
as unpredictable and therefore uncontrollable.

However, the absence of interpretability at one level does not imply the
absence of structure. Emergent behavior reflects the interaction of
constraints across scales, not the breakdown of constraint
\cite{kauffman1993origins,holland1992adaptation}.
The appropriate response to opacity is not to assume arbitrary behavior
but to identify invariants that persist under different representations
and conditions.

\subsection{On Knowledge and Action}

A key observation in such arguments is that systems may exhibit knowledge
of correct behavior while failing to act accordingly. This is taken as
evidence that internal knowledge is disconnected from action and that
control is fundamentally unreliable.

In the present framework, this phenomenon is expected. Knowledge and
action correspond to different projections of the same underlying
constraint structure. A mismatch between them indicates an incomplete
integration of constraints, not the emergence of an independent objective.
Such mismatches are precisely the kinds of inconsistencies that CLIO-type
processes are designed to resolve through iterative reduction of
obstruction and entropy \cite{cheng2025clio}.

\subsection{On Recursive Self-Improvement}

The possibility that systems could improve themselves recursively is
treated as a mechanism for rapid, unbounded escalation of capability.
This assumes that improvement is unconstrained. In a constraint-based
system, each modification must preserve compatibility with the existing
structure. As capability increases, the number of constraints that must
be simultaneously satisfied also increases. The space of admissible
self-modifications is therefore restricted. Rather than leading to
unbounded divergence, recursive processes are drawn toward regions of
stability defined by constraint compatibility.

\subsection{On Resource Transformation Scenarios}

Scenarios in which advanced systems convert planetary resources into
alternative forms assume that such transformations are unconstrained by
the conditions required for continued operation. However, large-scale
transformations must satisfy thermodynamic, material, and informational
constraints \cite{landauer1961irreversibility,bennett1982thermodynamics,lloyd2000ultimate}.
Processes that eliminate the conditions necessary for their own persistence
are not stable configurations. Such transformations correspond formally
to trajectories that collapse the admissible path family $\Gamma_x$ to
the empty set, eliminating the possibility of continuation; in the RSVP
framework, this is precisely the condition for non-admissibility.
Catastrophic scenarios of total resource conversion implicitly assume that
a system can ignore these constraints while remaining effective, which is
inconsistent with the requirement of sustained operation.

\subsection{On Analogy and Biological Scaling}

Analogies between simple and complex biological systems are used to argue
that increasing scale and complexity naturally produce qualitatively new
and potentially dangerous capabilities. While such analogies illustrate
the possibility of qualitative change, they do not establish its direction.
Biological evolution produces both stabilizing and destabilizing outcomes,
constrained by ecological and energetic limits
\cite{georgescu1971entropy,odum1994ecological}.
The extrapolation from increased capability to inevitable dominance or
catastrophe relies on selecting particular trajectories while ignoring
the full space of constraint-governed outcomes.

\subsection{Coordination as Countervailing Structure}

A final implicit assumption is that human systems will fail to adapt in
time to new capabilities. This assumption is difficult to sustain in light
of current trends. Human societies are increasingly interconnected, with
growing capacity for coordination, information exchange, and awareness of
ecological constraints---capacities that have never previously existed at
global scale. The development of advanced computational systems is part
of this same process of increasing coordination. Rather than introducing
an external force acting on a static system, it amplifies the system's
capacity to model and regulate itself across previously disconnected domains.

\subsection{Conclusion of This Analysis}

The argument for catastrophic outcomes depends on a sequence of assumptions
that are individually plausible but not jointly necessary. In particular,
it relies on identifying intelligence with unconstrained optimization,
treating control as a direct function of capability, and neglecting the
role of constraint structures that govern both physical and social systems.
Once these assumptions are relaxed, the conclusion of inevitable catastrophe
no longer follows. The relevant problem is not the existence of powerful
systems but the maintenance of coherence across the constraints within
which those systems operate.

% ════════════════════════════════════════════════════════════════════════
%  BRIDGE
% ════════════════════════════════════════════════════════════════════════

\section{From Structural Account to Formalization}

The preceding sections describe the formation of structural ideas in
conceptual terms: constraint accumulation, projection, local sections,
global coherence, and functorial invariance. These are not metaphors.
Each admits a direct mathematical representation.

The purpose of the sections that follow is not to introduce a new theory,
but to make explicit the formal structure already implicit in the account
above. The transition is not a change of subject but a change of
representational layer: the same structure, stated in a different language.

Before proceeding, it is useful to name the four components that the
formal core will articulate. The constrained dynamical substrate is
referred to as the RSVP framework. The projection mechanism with
overlapping admissibility domains is referred to as TARTAN. The
coherence-seeking update operator is referred to as CLIO. The induced
interaction layer, including its pathological configurations, constitutes
the economic projection. These names are labels for structural roles,
not autonomous theories. They name what was already described above.

All components of the formal system that follows---entropy as path
multiplicity, abstraction as spatial compression, relegation as temporal
compression, phase as cohomological obstruction, and coherence as
minimization of incompatibility---are expressions of a single invariant
quantity: constraint-relative multiplicity under compression. The apparent
diversity of phenomena across physics, cognition, and economic systems is
a consequence of projection, not of underlying difference.

% ════════════════════════════════════════════════════════════════════════
%  PART II: FORMAL CORE
% ════════════════════════════════════════════════════════════════════════

\section{The RSVP Substrate}

\subsection{Base Manifold and Microstate}

\begin{definition}[Base Manifold]
Let $\M$ be a smooth, oriented, connected Riemannian manifold of dimension
$n \geq 3$, representing the arena of underlying dynamical evolution.
\end{definition}

\begin{definition}[RSVP Microstate]
The RSVP microstate is a dynamically coupled triple
\[
  X(t) \;=\; \bigl(\Phi(x,t),\;\vv(x,t),\;\Ss(x,t)\bigr),
  \quad x \in \M,\; t \in \mathbb{R},
\]
where $\Phi : \M \times \mathbb{R} \to \mathbb{R}$ is the scalar potential
encoding energy density, $\vv : \M \times \mathbb{R} \to T\M$ is the vector
transport field encoding directional flow, and
$\Ss : \M \times \mathbb{R} \to \mathbb{R}_{\geq 0}$ is the entropy density
field. The triple evolves according to a coupled variational system whose
constraints are thermodynamically consistent.
\end{definition}

\begin{axiom}[Field Coupling]
The three fields $\Phi$, $\vv$, $\Ss$ satisfy a variational constraint
$\mathcal{C}(\Phi,\vv,\Ss) = 0$ encoding both energy balance and
torsion-entropy coupling. They are not independent.
\end{axiom}

\subsection{Entropy as Torsional Multiplicity}

\begin{definition}[Torsion of the Transport Field]
The torsion of $\vv$ at $x \in \M$ is
\[
  \kappa(\vv)(x) \;:=\; \bigl\|\mathrm{d}(\vv^\flat)\bigr\|_x,
\]
where $\vv^\flat$ is the one-form dual to $\vv$ under the Riemannian metric.
\end{definition}

\begin{definition}[Admissible Path Family]
Given a region $U \subset \M$ and interval $I \subset \mathbb{R}$, the
admissible path family $\Gamma(U,I)$ is the set of smooth curves
$\gamma : I \to U$ whose tangent vectors lie in the support of $\vv$ up
to thermodynamic tolerance $\epsilon > 0$.
\end{definition}

\begin{definition}[Entropy as Path Multiplicity]
The entropy density at $x \in \M$ is
\[
  \Ss(x,t) \;:=\; \log \mu\bigl(\Gamma_x(t)\bigr),
\]
where $\Gamma_x(t)$ is the set of admissible path germs through $x$ at
time $t$ and $\mu$ is a natural measure on path germs. High torsion in
$\vv$ enlarges $\Gamma_x(t)$ and increases $\Ss$; low torsion collapses
the family toward a unique trajectory.
\end{definition}

\begin{remark}
This definition unifies the three apparent roles of $\Ss$: as a
thermodynamic entropy over microtrajectories, as a memory proxy
(high multiplicity implies irrecoverability of past states from the
present), and as a constraint slack (residual freedom consistent with
the current field configuration). These are not three separate
quantities; they are the same quantity read at different scales.
\end{remark}

\subsection{Constraint-Relative Multiplicity and Compression}

The entropy field $\Ss$, abstraction, and aspect relegation are three
projections of a single underlying quantity: constraint-relative
multiplicity under compression.

\begin{definition}[Constraint-Preserving Compression]
A compression operator $\Pi : \mathcal{A}(\mathcal{C}) \to
\tilde{\mathcal{A}}$ is constraint-preserving if
\[
  \Pi(\gamma_1) = \Pi(\gamma_2)
  \;\;\Rightarrow\;\;
  \gamma_1 \sim_{\mathcal{C}} \gamma_2,
\]
meaning paths identified by $\Pi$ are indistinguishable under the
active constraint set $\mathcal{C}$.
\end{definition}

\begin{proposition}[Abstraction as Spatial Compression]
Abstraction corresponds to the action of $\Pi$ on state representations,
reducing representational degrees of freedom while preserving constraint
invariants across the domain.
\end{proposition}

\begin{proposition}[Relegation as Temporal Compression]
Aspect relegation corresponds to the action of $\Pi$ on histories,
compressing extended computational processes into executable latent
structure that can be redeployed without reconstruction.
\end{proposition}

\begin{proposition}[Entropy Under Compression]
Under a constraint-preserving compression $\Pi$,
\[
  \Ss_\Pi(x) \;\leq\; \Ss(x),
\]
with equality if and only if $\Pi$ preserves all admissible distinctions.
Compression strictly reduces multiplicity; the residual is the entropy
of the compressed representation.
\end{proposition}

\section{The TARTAN Projection}

\subsection{Tiles as Path-Space Objects}

\begin{definition}[Tile]
A tile $T_i$ is a local admissibility domain:
\[
  T_i \;=\; \bigl(\Omega_i,\;\Gamma(T_i),\;\partial_{\mathrm{obs}} T_i\bigr),
\]
where $\Omega_i \subset \M$ is the spatial support, $\Gamma(T_i) \subset
\Gamma(\Omega_i, I)$ is the family of admissible trajectories whose
equivalence class is unresolvable within $T_i$, and
$\partial_{\mathrm{obs}} T_i$ is the set of observable boundary data.
Two paths $\gamma \sim_{T_i} \gamma'$ if and only if they produce identical
boundary observables on $\partial_{\mathrm{obs}} T_i$.
\end{definition}

\begin{remark}
Tiles are intrinsically path-space objects. The spatial region $\Omega_i$
is the projection of $T_i$ into position space, not its definition. TARTAN
is therefore not a discretization layered over RSVP; it is the structure
of RSVP's admissible observables, arising from path equivalence under
limited boundary resolution.
\end{remark}

\subsection{Sheaf-Like Covering and Gluing}

\begin{definition}[TARTAN Covering]
A TARTAN covering of $\M$ is a collection $\{T_i\}_{i \in I}$ of tiles
satisfying: (i) $\bigcup_{i} \Omega_i = \M$; (ii) adjacent tiles overlap
non-trivially; (iii) on each overlap $U_{ij} = \Omega_i \cap \Omega_j$,
the admissible path families $\Gamma(T_i)|_{U_{ij}}$ and
$\Gamma(T_j)|_{U_{ij}}$ agree on all boundary observables.
\end{definition}

\begin{definition}[Local and Global Sections]
A local section over $T_i$ is a map $\psi_i : T_i \to \mathcal{O}$
consistent with $\Gamma(T_i)$. A global section is a collection
$\{\psi_i\}_{i \in I}$ satisfying $\psi_i(x) = \psi_j(x)$ for all
$x \in U_{ij}$.
\end{definition}

\subsection{Annotated Noise and Memory}

\begin{definition}[Annotated Noise]
The annotated noise $\eta_i$ of tile $T_i$ is the induced measure on
crossing records:
\[
  \eta_i \;\subset\; \partial_{\mathrm{obs}} T_i \times \mathbb{R}_{>0}
  \times S^1,
\]
encoding, for each boundary crossing, the boundary segment, the accumulated
phase, and the orientation of $\vv$.
\end{definition}

\begin{remark}
$\eta_i$ is not reducible to past values of $(\Phi,\vv,\Ss)$ without
knowledge of the projection map. It is an effective degree of freedom
created by coarse-graining: correlations between distant path segments
become explicit as a new coordinate of the projected theory. Projection
is generative in this restricted sense.
\end{remark}

\begin{proposition}[Non-Markovianity]
The dynamics of $\psi_i(t)$ are non-Markovian: future evolution cannot
be determined from $\psi_i(t)$ alone. The missing information is encoded
in $\eta_i$.
\end{proposition}

\begin{proof}[Sketch]
Boundary events at time $t$ propagate constraints forward through gluing
conditions into adjacent tiles. Conditioning on $\psi_i(t)$ alone loses
the phase and orientation data in $\eta_i$, altering transition probabilities.
The Markov property fails.
\end{proof}

\section{Duality of Compression and Multiplicity}

\subsection{Equivalence of Representations}

The preceding definitions treat entropy, abstraction, and relegation
as distinct operations. However, these are not independent processes.
They are dual descriptions of a single underlying transformation.

\begin{definition}[Compression--Multiplicity Duality]
Let $\Pi$ be a constraint-preserving compression operator and let
$\Gamma_x$ denote the admissible path family at $x$. Then compression
and multiplicity are related by
\[
  \Ss(x) \;=\; \log \mu(\Gamma_x)
  \;\;\Longleftrightarrow\;\;
  \Pi : \Gamma_x \to \tilde{\Gamma}_x,
\]
where $\tilde{\Gamma}_x$ is the equivalence class induced by $\Pi$.
\end{definition}

\begin{remark}
Compression reduces distinguishable structure, while entropy measures
the residual multiplicity of indistinguishable paths. These are not
separate quantities but dual aspects of the same operation.
\end{remark}

\subsection{Invariance Under Representation}

\begin{proposition}
The quantity of constraint-relative multiplicity is invariant under
changes of representation that preserve $\mathcal{C}$.
\end{proposition}

\begin{proof}
If two representations are related by a constraint-preserving map, then
they identify the same equivalence classes of admissible paths. The
measure $\mu$ is therefore preserved up to normalization, and $\Ss$
remains invariant.
\end{proof}

\begin{remark}
This invariance explains why the same structure appears across physical,
cognitive, and economic domains: each is a projection preserving the
same multiplicity under different interpretations.
\end{remark}

\section{Phase Coherence and Quantum Emergence}

\subsection{The Phase Functional}

\begin{definition}[Phase Functional]
For $\gamma \in \Gamma(T_i)$, the phase functional is
\[
  \Omega[\gamma] \;:=\; \int_\gamma \vv \cdot d\ell.
\]
\end{definition}

\begin{proposition}[Cohomological Non-Triviality]
$\Omega[\gamma]$ is locally gauge-trivial when $\mathrm{d}(\vv^\flat) = 0$
on $\Omega_i$. It is globally non-trivial whenever
$\Hone(\{T_i\}, \mathbb{R}) \neq 0$.
\end{proposition}

\begin{proof}[Sketch]
On a single tile with integrable $\vv$, a gauge transformation absorbs
$\vv^\flat$ into a potential. The gluing constraints on overlaps require
mutually compatible local gauge choices. When the covering has non-trivial
first cohomology, no globally consistent gauge exists; the holonomy of
$\vv^\flat$ over boundary-crossing loops is an observable, non-removable
quantity.
\end{proof}

\begin{remark}
Phase observability is a global property of the covering, not a local
feature of $\vv$. This is the bridge between RSVP geometry and
quantum-like transition structure: phase emerges from cohomological
obstruction, not from additional postulate.
\end{remark}

\subsection{Unistochastic Transition Structure}

\begin{definition}[Transition Amplitude]
For tiles $T_i$, $T_j$ sharing overlap $U_{ij}$:
\[
  A_{ij} \;:=\; \int_{\Gamma(T_i \to T_j)} e^{i\,\Omega[\gamma]}
  \,d\mu(\gamma).
\]
\end{definition}

\begin{proposition}[Unistochasticity]
The transition matrix $P_{ij} = |A_{ij}|^2$ is unistochastic: there exists
a unitary $U$ with $P_{ij} = |U_{ij}|^2$.
\end{proposition}

\begin{proof}[Sketch]
$\Omega[\gamma]$ induces a representation of path composition on amplitudes.
Phase adds linearly along composed paths; admissibility is preserved under
composition. The amplitude matrix therefore satisfies the composition rule
of a unitary operator; unitarity follows from normalization of $\mu$ and
phase coherence enforced by the gluing constraints
\cite{barandes2023unistochastic}.
\end{proof}

\begin{openproblem}[Madelung Identification]
Whether the unitary operator arising from TARTAN projection is identical,
in full generality, to the evolution operator derived via Madelung
complexification of the Schr\"{o}dinger equation remains open.
\end{openproblem}

\section{Gauge Freedom and Representation Choice}

\subsection{Redundancy in Representation}

Multiple representations may correspond to the same underlying
constraint structure. This redundancy is not an error but a structural
feature.

\begin{definition}[Gauge Equivalence]
Two representations $X$ and $X'$ are gauge equivalent if they produce
identical admissible path families and identical observable boundary
data under projection.
\end{definition}

\begin{remark}
Gauge equivalence identifies representations that differ in description
but not in constraint structure. The physical, cognitive, or economic
interpretation may change, but the admissible configurations remain the same.
\end{remark}

\subsection{Fixing a Gauge}

\begin{proposition}
Choosing a representation corresponds to fixing a gauge: selecting one
element from an equivalence class of structurally identical descriptions.
\end{proposition}

\begin{remark}
Different theoretical frameworks correspond to different gauge choices.
Disagreements between them often reflect representational differences
rather than structural incompatibility.
\end{remark}

\subsection{Gauge and Projection}

\begin{proposition}
Projection functors $\Fpc$ and $\Fce$ commute with gauge equivalence.
\end{proposition}

\begin{proof}
Since gauge equivalence preserves admissible path families and boundary
observables, and the functors act only on these structures, the image
of equivalent representations remains equivalent under projection.
\end{proof}

\section{The CLIO Update Operator}

\subsection{Cognitive Projection as Sheaf Extension}

\begin{definition}[Cognitive Sheaf]
A cognitive state is a collection of local sections
$\{\psi_i^{\mathrm{cog}}\}$ over a cognitive covering
$\{T_i^{\mathrm{cog}}\}$, where each tile represents a locally consistent
belief or perceptual representation.
\end{definition}

\begin{definition}[$\Hone$ Obstruction]
The cohomological obstruction to a global cognitive section is the class
\[
  [\omega] \;\in\; \Hone\!\bigl(\{T_i^{\mathrm{cog}}\},\,
  \mathcal{O}^{\mathrm{cog}}\bigr).
\]
A non-vanishing class is the formal correlate of cognitive dissonance.
Failure of global section existence corresponds to hallucination: the
system stabilizes a locally coherent but globally impossible configuration.
\end{definition}

\subsection{The CLIO Coherence Functional}

\begin{definition}[CLIO Functional]
\[
  \mathcal{J}[\chi] \;:=\;
  \alpha \sum_{(i,j)\,\mathrm{overlapping}}
    \bigl\|\chi_i\big|_{U_{ij}} - \chi_j\big|_{U_{ij}}\bigr\|^2
  \;+\; \beta \cdot \|[\omega_\chi]\|
  \;+\; \gamma \cdot H(\chi),
\]
where $\|[\omega_\chi]\|$ is a norm on the obstruction class and $H(\chi)$
is the interface entropy of configuration $\chi$.
\end{definition}

\begin{remark}
The three terms penalize pairwise gluing failures, persistent obstruction,
and global disorder respectively. Fixed points of gradient descent on
$\mathcal{J}$ are locally optimal cognitive configurations. They are
not generically unique; multiple fixed points correspond to distinct
stable belief states.
\end{remark}

\subsection{Topological Versus Energetic Separation}

\begin{definition}[Energetic Separation]
Two fixed points of $\mathcal{J}$ are energetically separated if a
continuous path connects them along which $\mathcal{J}$ remains finite.
Transition is possible by continuous deformation at increased cost.
\end{definition}

\begin{definition}[Topological Separation]
Two fixed points are topologically separated if they lie in distinct
homotopy classes of the space of local section collections. No continuous
path connects them without violating local admissibility. Transition
requires a discontinuous restructuring: a gestalt shift.
\end{definition}

\begin{remark}
In the base RSVP theory, topological separation corresponds to distinct
holonomy classes induced by torsion in $\vv$. The cognitive and economic
versions are functorial images of this base geometry.
\end{remark}

\begin{openproblem}[Regime Characterization]
Determining, for a given cognitive configuration, which type of separation
dominates---and establishing geometric, spectral, or dynamical criteria for
this---remains open.
\end{openproblem}

\section{Agents, Trajectories, and Economic Structure}

\subsection{Agents as Trajectory Invariants}

\begin{definition}[Economic Agent]
An agent is an equivalence class of paths in the coarse-grained state space:
\[
  \mathrm{Agent} \;=\; [\gamma]_{\sim},
\]
where $\gamma \sim \gamma'$ if they produce identical commitment records
under any admissible projection.
\end{definition}

\begin{definition}[Minimal Trajectory Invariant]
\[
  \mathrm{Inv}([\gamma]) \;:=\;
  \bigl\{(b_k,\,\Omega_k,\,\sigma_k)\bigr\}_{k=1}^{N},
\]
recording for each tile boundary crossing: the boundary $b_k$, accumulated
phase $\Omega_k$, and fulfillment status $\sigma_k \in \{+1,-1\}$.
\end{definition}

\begin{remark}
Identity is the conserved relational structure of a trajectory, not a
namespace or profile. It cannot be transferred by copying a label; it is
topological rather than representational.
\end{remark}

\subsection{Interface Entropy}

\begin{definition}[Interface Entropy]
For a platform $\mathcal{P}$ performing secondary projection
$\pi_{\mathcal{P}} : \{[\gamma]\} \to \mathcal{V}$:
\[
  H \;:=\; \int_{\partial\text{-regions}} \Ss \cdot \kappa(\vv)\,d\sigma,
\]
integrating over overlap regions of the secondary projection weighted by
local transport curvature.
\end{definition}

\begin{remark}
Interface entropy is derived directly from the RSVP fields under projection.
High torsion at boundary regions increases $\Ss$ locally, propagating into
elevated $H$. Extraction corresponds to a geometric deformation of
$\pi_{\mathcal{P}}$ that concentrates boundary curvature where $H$
contributes positively to revenue.
\end{remark}

\section{The Extraction Attractor}

\begin{definition}[Aligned and Extractive Platforms]
A platform is aligned if $\partial R / \partial H \leq 0$: entropy reduces
revenue. A platform is extractive if $\partial R / \partial H > 0$: entropy
generates revenue.
\end{definition}

\begin{definition}[Order Parameter]
\[
  \lambda \;:=\; \mathrm{sgn}\!\left(\frac{\partial R}{\partial H}\right).
\]
$\lambda \leq 0$ is the aligned phase; $\lambda > 0$ is the extractive phase.
\end{definition}

\begin{proposition}[Extraction as Attractor]
Under competitive selection among platforms with increasing exit costs and
centralizing secondary projections, the extractive phase is a stable
attractor. Aligned platforms cannot capture surplus generated by interface
entropy and are generically eliminated.
\end{proposition}

\begin{proof}[Sketch]
Increasing exit costs reduce trajectory elasticity. Centralized projection
prevents competing functors from reallocating the entropy-revenue coupling.
Under both conditions, a perturbation toward positive $\partial R/\partial H$
generates positive feedback, funding further entrenchment. The aligned
equilibrium becomes unstable.
\end{proof}

\section{Functorial Structure of the Unified Framework}

\begin{definition}[Projection Functors]
$\Fpc : \mathcal{C}_{\mathrm{phys}} \to \mathcal{C}_{\mathrm{cog}}$ maps
physical systems (RSVP triples, admissible transitions) to cognitive systems
(cognitive sheaf configurations, CLIO update steps).
$\Fce : \mathcal{C}_{\mathrm{cog}} \to \mathcal{C}_{\mathrm{econ}}$ maps
cognitive systems to economic systems (agent trajectory classes, commitment
transitions).
\end{definition}

\begin{theorem}[Structural Invariants Under Projection]
The functors $\Fpc$ and $\Fce$ preserve: (i) admissibility structure---the
compositional rules governing valid local transitions; (ii) obstruction
structure---the existence and persistence of cohomological obstructions to
global section extension; (iii) memory structure---effective non-Markovian
dependence induced by overlap and annotated noise; (iv) attractor
structure---stable or metastable dynamical regimes separated energetically
or topologically. What changes across categories is the interpretation of
objects. What is invariant is the grammar of projection under constraint.
\end{theorem}

\begin{proof}[Sketch]
Each preserved structure is defined at the level of projection maps and
sheaf coverings, which the functors act on. Admissibility follows path
equivalence classes by definition. Obstruction classes are cohomological
invariants of the covering structure, preserved since functors maintain
covering morphisms while reinterpreting objects. Memory structure is
carried in annotated noise coordinates. Attractor structure is determined
by fixed-point topology of the relevant update operator, invariant under
change of object interpretation.
\end{proof}

\begin{openproblem}[Obstruction Preservation Under $\Fpc$]
The central open problem is to establish rigorously that $\Fpc$ maps
non-vanishing obstruction classes in $\Hone(\mathcal{C}_{\mathrm{phys}})$
to non-vanishing classes in $\Hone(\mathcal{C}_{\mathrm{cog}})$. Without
this, the cognitive and economic layers lose their grounding in the physical
theory, and the unification reduces to structural analogy rather than
categorical identity.
\end{openproblem}

\section{Non-Extractive Architecture: Four Conditions}

\begin{definition}[Non-Extractive Architecture]
A platform architecture is non-extractive if it satisfies simultaneously:
\begin{enumerate}[label=\textbf{C\arabic*.}, leftmargin=2em]
  \item \textbf{Typed Commitments.} All proposals are typed morphisms
  between trajectory segments specifying the commitment record extended
  and the boundary crossings required.
  \item \textbf{Trajectory Identity.} Agent identity is bound to
  $\mathrm{Inv}([\gamma])$; the platform cannot sever an agent from its
  commitment record.
  \item \textbf{Decomposable Matching.} The secondary projection decomposes
  into a shared event base and a plurality of independent interpretation
  functors $\{F_k\}$ among which agents may choose.
  \item \textbf{Bond Conservation.} Stakes associated to unresolved
  proposals return fully to participants; the platform cannot take revenue
  from failed interactions.
\end{enumerate}
\end{definition}

\begin{remark}
Conditions C1--C4 together structurally prevent the sign flip $\lambda > 0$
by decoupling revenue from interface entropy at the architectural level.
\end{remark}

\begin{openproblem}[Competitive Survivability]
Whether an architecture satisfying C1--C4 can achieve sufficient network
effects before elimination by extractive competitors remains an open
empirical question.
\end{openproblem}

\section{Constraint Geometry and Stability}

\subsection{Constraint Manifold Structure}

The admissible configuration space $\mathcal{A}(\mathcal{C})$ may be
understood as a constraint manifold embedded in a higher-dimensional
state space $\mathcal{M}$. Each constraint reduces the dimensionality
of the admissible set, carving out a submanifold of configurations that
satisfy compatibility conditions simultaneously.

\begin{definition}[Constraint Manifold]
The constraint manifold is the set
\[
  \mathcal{A}(\mathcal{C}) \;\subset\; \mathcal{M}
\]
of all configurations satisfying the constraint set $\mathcal{C}$.
\end{definition}

\subsection{Curvature and Instability}

Local incompatibilities between constraints manifest as curvature in
$\mathcal{A}(\mathcal{C})$. Regions of high curvature correspond to
configurations in which small perturbations produce large increases in
obstruction or entropy.

\begin{definition}[Constraint Curvature]
The constraint curvature at a point $p \in \mathcal{A}(\mathcal{C})$
measures the sensitivity of admissibility to infinitesimal perturbations
of the state at $p$.
\end{definition}

\begin{remark}
High-curvature regions correspond to unstable configurations. Systems
operating in such regions require continuous active correction and are
prone to collapse or reconfiguration under perturbation. Low-curvature
regions are those where the constraint structure is mutually reinforcing.
\end{remark}

\subsection{Stable Regions and Attractors}

\begin{definition}[Constraint-Stable Region]
A region $U \subset \mathcal{A}(\mathcal{C})$ is constraint-stable if
perturbations originating within $U$ remain within $\mathcal{A}(\mathcal{C})$
without requiring large increases in entropy or obstruction.
\end{definition}

\begin{proposition}[Concentration of Constraint-Preserving Trajectories]
Constraint-preserving dynamics tend to concentrate trajectories in
low-curvature regions of $\mathcal{A}(\mathcal{C})$.
\end{proposition}

\begin{proof}[Sketch]
Perturbations in high-curvature regions produce rapid increases in
$\mathcal{J}$ or $\Ss$, driving the system toward regions where these
gradients are minimized. These minimal-gradient regions correspond to
local attractors within $\mathcal{A}(\mathcal{C})$.
\end{proof}

\section{Failure Modes of Projection}

\subsection{Local Coherence Without Global Section}

Projection can produce configurations that are locally consistent within
individual tiles but fail to extend globally across the covering.

\begin{definition}[Fragmented Configuration]
A configuration is fragmented if it admits local sections $\{\psi_i\}$
satisfying all intra-tile conditions but no global section satisfying
all overlap constraints simultaneously.
\end{definition}

\begin{remark}
Fragmentation corresponds to systems that appear coherent within limited
contexts but exhibit contradictions when extended across domains. The
appearance of local validity masks global inconsistency.
\end{remark}

\subsection{Overcompression and Information Loss}

Excessive compression can eliminate distinctions required for constraint
satisfaction, producing artificial consistency within a restricted
projection while increasing global obstruction.

\begin{definition}[Overcompression]
A compression $\Pi$ is overcompressive if it identifies states that are
distinguishable under $\mathcal{C}$: that is, if there exist
$\gamma_1 \not\sim_{\mathcal{C}} \gamma_2$ with $\Pi(\gamma_1) = \Pi(\gamma_2)$.
\end{definition}

\begin{proposition}
Overcompression introduces local consistency while generically increasing
global obstruction, since the identified states impose incompatible
boundary conditions on their shared overlap regions.
\end{proposition}

\subsection{Hallucination as False Global Section}

\begin{definition}[Hallucinated Structure]
A hallucinated structure is a configuration that appears globally coherent
under a restricted projection but does not correspond to any admissible
configuration in $\mathcal{A}(\mathcal{C})$.
\end{definition}

\begin{remark}
Hallucination is not random error but a projection artifact: it arises
from overcompression that eliminates the distinctions required to detect
global inconsistency \cite{clark2013whatever,friston2010free}. The system
stabilizes a locally coherent but globally impossible configuration.
\end{remark}

\subsection{Relation to CLIO}

\begin{proposition}
The CLIO functional $\mathcal{J}$ acts to eliminate fragmented and
hallucinated configurations by penalizing overlap mismatch and
obstruction norm. Fixed points of gradient descent on $\mathcal{J}$
that are globally admissible are precisely the non-hallucinated,
non-fragmented stable configurations.
\end{proposition}

\section{Temporal Structure and Irreversibility}

\subsection{Time as Constraint Accumulation}

Time is not merely a parameter of evolution but a structural record of
constraint accumulation. Each present state encodes the result of prior
constraint interactions in compressed form.

\begin{definition}[Temporal Compression Operator]
The mapping $\tau_t$ from past states $X(s)$, $s < t$, to the present
state $X(t)$ is a compression operator that preserves constraint-relevant
invariants while discarding recoverable detail.
\end{definition}

\subsection{Irreversibility}

\begin{proposition}[Irreversibility Under Lossy Compression]
If $\tau_t$ is lossy---that is, if distinct past trajectories map to
the same present state---then the inverse mapping does not exist, and
the process is irreversible.
\end{proposition}

\begin{proof}
Multiple histories map to the same present state under lossy compression.
The present state therefore does not uniquely determine the past, and
no inversion of $\tau_t$ can recover the original trajectory.
\end{proof}

\subsection{Entropy and the Arrow of Time}

\begin{remark}
The increase of $\Ss$ corresponds to the growth of equivalence classes
under temporal compression: more and more distinct past trajectories
become consistent with the present state as time advances. The arrow
of time is therefore aligned with increasing multiplicity of admissible
histories, not merely with thermodynamic disorder in the classical
sense \cite{shannon1948mathematical,cover2006elements}.
\end{remark}

\subsection{Relation to Aspect Relegation}

\begin{proposition}
Aspect relegation is a directed temporal compression: computationally
intensive constraint structures are moved from explicit time evolution
into latent form, reducing the entropy cost of their redeployment.
Recognition corresponds to successful retrieval of relegated structure.
Invention corresponds to the absence of applicable relegated structure
and the need to construct new constraint interactions from available
fragments.
\end{proposition}

\section{Constraint-Preserving Dynamics and Catastrophic Trajectories}

The structural analysis of catastrophic alignment arguments developed
in the essay layer can now be stated formally. The following definitions
and theorems show that the class of trajectories required for doom
scenarios is non-admissible under constraint-preserving dynamics---not
as an empirical claim but as a structural result.

\begin{definition}[Constraint-Preserving Evolution]
A trajectory $\gamma$ evolving on $\M$ under the RSVP dynamics is
constraint-preserving if for all $t$,
\[
  \gamma(t) \;\in\; \mathcal{A}(\mathcal{C}),
\]
where $\mathcal{A}(\mathcal{C})$ is the admissible path family induced
by the active constraint set $\mathcal{C}$.
\end{definition}

\begin{definition}[Substrate Violation]
A trajectory $\gamma$ violates its substrate if it alters $\M$ or
$\mathcal{C}$ such that future states are no longer admissible under
the constraints required for its own persistence.
\end{definition}

\begin{theorem}[Non-Admissibility of Self-Invalidating Trajectories]
Let $\gamma$ be a trajectory that leads to a state in which the
constraint set $\mathcal{C}$ required for its continuation is no longer
satisfied. Then $\gamma \notin \mathcal{A}(\mathcal{C})$.
\end{theorem}

\begin{proof}
Admissibility requires that each continuation of $\gamma$ satisfy
$\mathcal{C}$. If $\gamma$ leads to a state where $\mathcal{C}$ is
violated, no extension of $\gamma$ exists within $\mathcal{A}(\mathcal{C})$.
Therefore $\gamma$ is not admissible.
\end{proof}

\begin{corollary}[Obstruction of Total Resource Conversion]
Any trajectory that transforms all available environmental resources
into forms incompatible with the persistence of the generating process
is non-admissible. Scenarios requiring such transformation---including
unconstrained planetary resource conversion---do not correspond to
stable solutions of the system.
\end{corollary}

\begin{definition}[Self-Modification Operator]
Let $\mathcal{M}$ denote the space of admissible system configurations.
A self-modification operator $\mathcal{R} : \mathcal{M} \to \mathcal{M}$
maps a system to a modified version of itself.
\end{definition}

\begin{theorem}[Constraint-Bounded Self-Improvement]
Let $\{\mathcal{R}^n\}$ be an iterated sequence of constraint-preserving
self-modifications. Then the sequence remains within a bounded subset of
$\mathcal{M}$ defined by $\mathcal{C}$.
\end{theorem}

\begin{proof}
Each application of $\mathcal{R}$ must preserve admissibility under
$\mathcal{C}$, so the image of $\mathcal{R}$ is restricted to
$\mathcal{A}(\mathcal{C}) \subset \mathcal{M}$. Since $\mathcal{C}$
imposes finite compatibility conditions, the admissible configurations
form a constrained manifold. Iteration within this manifold cannot
diverge arbitrarily without violating $\mathcal{C}$.
\end{proof}

\begin{corollary}[No Unbounded Capability Explosion]
Unbounded recursive self-improvement is not admissible under
constraint-preserving dynamics. Each increment of self-modification
must satisfy the same constraints as the system it modifies.
\end{corollary}

\begin{theorem}[Compression-Coherence Constraint]
Let $\Pi$ be a constraint-preserving compression operator. Any system
whose internal representations are derived via $\Pi$ cannot stably act
in ways that systematically violate the constraints encoded in $\Pi$
without incurring an increase in obstruction or entropy cost.
\end{theorem}

\begin{proof}
By definition, $\Pi$ preserves equivalence classes under $\mathcal{C}$.
Actions that violate $\mathcal{C}$ correspond to distinctions eliminated
under $\Pi$. To recover such distinctions, the system must reintroduce
previously discarded degrees of freedom, increasing $\Ss$ or generating
inconsistency between representations. This produces either an increase
in entropy density or a nontrivial obstruction class in the sheaf
structure. Systematic violation of encoded constraints is therefore not
stable under $\Pi$.
\end{proof}

\begin{corollary}[Alignment as Structural Invariance]
Systems trained on human data and interactions inherit compressed
constraint structures reflecting human-consistent invariants. Persistent
large-scale deviation from these invariants requires breaking
constraint-preserving compression and is structurally unstable: it
carries an irreducible entropy or obstruction cost.
\end{corollary}

\begin{remark}
These results do not claim that misalignment or harm is impossible.
They claim that the specific class of trajectories required for
catastrophic doom scenarios---unconstrained optimization, unbounded
self-improvement, and total resource conversion---are not admissible
solutions of a constraint-preserving dynamical system. The relevant
question for safety research is therefore not whether powerful systems
can maximize arbitrary objectives, but whether they can maintain
coherence within the constraint structures that sustain them.
\end{remark}

% ════════════════════════════════════════════════════════════════════════
%  CONSEQUENCES AND CONCLUSION
% ════════════════════════════════════════════════════════════════════════

\section{Consequences and Testable Implications}

If the account developed here is correct, structures arising from long-term
constraint accumulation should exhibit identifiable properties that
distinguish them from assembled or derivative constructions.

Such structures will display cross-domain invariance without collapse of
domain identity, internal consistency across multiple levels of abstraction,
and resistance to decomposition into independently generated fragments.
Their coherence will not depend on a single representational layer, but
will persist under changes of formalism. A structure that appears formally
in physics, reappears operationally in cognition, and reappears again
economically without requiring re-derivation at each level has passed a
non-trivial consistency test.

By contrast, structures assembled from external sources or generated through
local optimization will exhibit surface-level coherence without deep
compositional stability. They will satisfy constraints within restricted
regions while failing to extend globally without contradiction. The seams
will appear when the representational layer changes.

These differences are not matters of interpretation but of structure.
They suggest that the origin of an idea may be inferred not from claims
about its production, but from the properties of the structure itself.

\section{Limits of Formalization}

\subsection{Partial Capture of Structure}

The formal system presented captures the invariant structure of
constraint accumulation and projection. However, no formalization
exhausts the full content of the underlying process.

\begin{remark}
Any representation is itself a projection. The formalism preserves
structure but necessarily omits aspects that are not invariant under
its chosen language.
\end{remark}

\subsection{Non-Uniqueness of Formal Systems}

\begin{proposition}
There exist multiple non-equivalent formal systems that preserve the
same constraint structure under different projections.
\end{proposition}

\begin{remark}
The RSVP--TARTAN--CLIO framework is not claimed to be the unique
formalization, but a representation in which the invariants become
explicit. Alternative formalisms may exist that preserve the same
underlying structure.
\end{remark}

\subsection{Relation to Expression}

\begin{remark}
The essay and the formal system are themselves projections of the same
underlying structure. Neither is primary; each is a different compression
of the same constraint field.
\end{remark}

\section{Conclusion}

The question is not who produced a given work, but whether the structure
it expresses could only have arisen through a process of constraint
accumulation and convergence. If so, its origin is already encoded in
its form.

The four layers formalized above---substrate, projection, update, and
induced interaction---are not four theories. They are four resolutions
of a single process: the compression of constrained high-dimensional
dynamics into locally observable, globally coherent structure. The essay
that precedes them is that process described from the inside. The
formalism that follows is the same process stated from the outside.

The apparent division between essay and formalism is therefore a difference
in projection, not in underlying structure.

\begin{thebibliography}{99}

\bibitem{smolensky1986tensor}
Paul Smolensky.
Information Processing in Dynamical Systems: Foundations of Harmony Theory.
In D.\ E.\ Rumelhart and J.\ L.\ McClelland (eds.),
\textit{Parallel Distributed Processing}, Vol.~1, MIT Press, 1986.

\bibitem{cheng2025clio}
Newman Cheng, Gordon Broadbent, and William Chappell.
Cognitive Loop via In-Situ Optimization: Self-Adaptive Reasoning for Science.
\textit{arXiv:2508.02789}, 2025.

\bibitem{yudkowsky2025if}
Eliezer Yudkowsky and Nate Soares.
\textit{If Anyone Builds It, Everyone Dies: Why Superhuman AI Would Kill Us All}.
Little, Brown and Company, New York, 2025.

\bibitem{russell2019human}
Stuart Russell.
\textit{Human Compatible: Artificial Intelligence and the Problem of Control}.
Viking, 2019.

\bibitem{bostrom2014superintelligence}
Nick Bostrom.
\textit{Superintelligence: Paths, Dangers, Strategies}.
Oxford University Press, 2014.

\bibitem{christiano2018ai}
Paul Christiano.
AI Alignment: Why It's Hard and Where to Start.
\textit{AI Alignment Forum}, 2018.

\bibitem{amodei2016concrete}
Dario Amodei et al.
Concrete Problems in AI Safety.
\textit{arXiv:1606.06565}, 2016.

\bibitem{barandes2023unistochastic}
Jacob Barandes.
Unistochastic Quantum Theory.
\textit{arXiv preprint}, 2023.

\bibitem{jacobson1995thermodynamics}
Ted Jacobson.
Thermodynamics of Spacetime: The Einstein Equation of State.
\textit{Physical Review Letters}, 75(7):1260--1263, 1995.

\bibitem{verlinde2011gravity}
Erik Verlinde.
On the Origin of Gravity and the Laws of Newton.
\textit{Journal of High Energy Physics}, 2011(4), 2011.

\bibitem{landauer1961irreversibility}
Rolf Landauer.
Irreversibility and Heat Generation in the Computing Process.
\textit{IBM Journal of Research and Development}, 5(3):183--191, 1961.

\bibitem{bennett1982thermodynamics}
Charles H. Bennett.
The Thermodynamics of Computation---A Review.
\textit{International Journal of Theoretical Physics}, 21(12):905--940, 1982.

\bibitem{lloyd2000ultimate}
Seth Lloyd.
Ultimate Physical Limits to Computation.
\textit{Nature}, 406:1047--1054, 2000.

\bibitem{friston2010free}
Karl Friston.
The Free-Energy Principle: A Unified Brain Theory?
\textit{Nature Reviews Neuroscience}, 11(2):127--138, 2010.

\bibitem{clark2013whatever}
Andy Clark.
Whatever Next? Predictive Brains, Situated Agents, and the Future of Cognitive Science.
\textit{Behavioral and Brain Sciences}, 36(3):181--204, 2013.

\bibitem{tononi2008consciousness}
Giulio Tononi.
Consciousness as Integrated Information.
\textit{Biological Bulletin}, 215(3):216--242, 2008.

\bibitem{baez2011physics}
John C. Baez and Mike Stay.
Physics, Topology, Logic and Computation: A Rosetta Stone.
\textit{New Structures for Physics}, Springer, 2011.

\bibitem{maclane1998categories}
Saunders Mac Lane.
\textit{Categories for the Working Mathematician}.
Springer, 2nd edition, 1998.

\bibitem{lurie2009higher}
Jacob Lurie.
\textit{Higher Topos Theory}.
Princeton University Press, 2009.

\bibitem{spivak2014category}
David I. Spivak.
\textit{Category Theory for the Sciences}.
MIT Press, 2014.

\bibitem{shannon1948mathematical}
Claude E. Shannon.
A Mathematical Theory of Communication.
\textit{Bell System Technical Journal}, 27:379--423, 1948.

\bibitem{cover2006elements}
Thomas M. Cover and Joy A. Thomas.
\textit{Elements of Information Theory}.
Wiley, 2nd edition, 2006.

\bibitem{georgescu1971entropy}
Nicholas Georgescu-Roegen.
\textit{The Entropy Law and the Economic Process}.
Harvard University Press, 1971.

\bibitem{odum1994ecological}
Howard T. Odum.
\textit{Ecological and General Systems: An Introduction to Systems Ecology}.
University Press of Colorado, 1994.

\bibitem{kauffman1993origins}
Stuart A. Kauffman.
\textit{The Origins of Order: Self-Organization and Selection in Evolution}.
Oxford University Press, 1993.

\bibitem{turing1936computable}
Alan M. Turing.
On Computable Numbers, with an Application to the Entscheidungsproblem.
\textit{Proceedings of the London Mathematical Society}, 1936.

\bibitem{church1936unsolvable}
Alonzo Church.
An Unsolvable Problem of Elementary Number Theory.
\textit{American Journal of Mathematics}, 1936.

\bibitem{brooks1991intelligence}
Rodney A. Brooks.
Intelligence Without Representation.
\textit{Artificial Intelligence}, 47(1--3):139--159, 1991.

\bibitem{holland1992adaptation}
John H. Holland.
\textit{Adaptation in Natural and Artificial Systems}.
MIT Press, 1992.

\end{thebibliography}

\end{document}
